\documentclass[DIV=14]{scrartcl}
\usepackage[utf8]{inputenc}
\usepackage[T1]{fontenc}
\usepackage[ngerman]{babel}
\usepackage{graphicx}
\graphicspath{{/home/luise/Dokumente/StAPF/}}
\usepackage[draft, markup=underlined]{changes}
\usepackage{csquotes}

\usepackage{ulem}
%\usepackage[dvipsnames]{xcolor}
\usepackage{paralist}
\usepackage{fixltx2e}
%\usepackage{ellipsis}
\usepackage[tracking=true]{microtype}

\usepackage{lmodern}                        % Ersatz fuer Computer Modern-Schriften
%\usepackage{hfoldsty}

%\usepackage{fourier}             % Schriftart
\usepackage[scaled=0.81]{helvet}     % Schriftart

\usepackage{url}
\usepackage{tocloft}             % Paket für Table of Contents

\usepackage{xcolor}
\definecolor{urlred}{HTML}{660000}

\usepackage{hyperref}
\hypersetup{
    colorlinks=true,    
    linkcolor=black,    % Farbe der internen Links (u.a. Table of Contents)
    urlcolor=black,    % Farbe der url-links
    citecolor=black} % Farbe der Literaturverzeichnis-Links

\usepackage{mdwlist} %Änderung Zeilenabstände itemize & enumerate

%\usepackage{draftwatermark} %Wasserzeichen ``Entwurf'' 
%\SetWatermarkText{Entwurf}

\parindent 0pt                 % Absatzeinrücken verhindern
\parskip 12pt                 % Absätze durch Lücke trennen

\setlength{\textheight}{23cm}
\usepackage{fancyhdr}
\pagestyle{fancy}
\fancyhead{} % clear all header fields
\cfoot{}
\lfoot{Zusammenkunft aller Physik-Fachschaften}
\rfoot{www.zapfev.de\\stapf@zapf.in}
\renewcommand{\headrulewidth}{0pt}
\renewcommand{\footrulewidth}{0.1pt}
\newcommand{\gen}{*innen}
\addto{\captionsngerman}{\renewcommand{\refname}{Quellen}}

%%%% Mit-TeXen Kommandoset
\usepackage[normalem]{ulem}
\usepackage{xcolor}
\newcommand{\replace}[2]{
    \sout{\textcolor{blue}{#1}}~\textcolor{blue}{#2}
}
\newcommand{\delete}[1]{
    \sout{#1}
}
\newcommand{\add}[1]{
    \textcolor{blue}{#1}
}
\newcommand{\okay}[1]{
    \textcolor{black}{#1}
}





\begin{document}

\hspace{0.87\textwidth}
\begin{minipage}{120pt}
\vspace{-1.8cm}
%\includegraphics[width=80pt]{logo.pdf}
\centering
\small Zusammenkunft aller Physik-Fachschaften
\end{minipage}

\begin{center}
\huge{Resolution der Zusammenkunft aller Physik Fachschaften}\vspace{.25\baselineskip}\\
\normalsize
\end{center}
\vspace{1cm}


\section*{Zur geplanten Novelle des Berliner Hochschulgesetzes (BerlHG)}

%\textbf{Antragsteller*innen:} Jenny (FUB), Hannah (HUB), Stephie (HUB), lobachevsky (FUB), scatty (RWTH)

%\textbf{Adressaten:}
%\begin{itemize}
%%\tightlist
%\item wissenschaftliche Sprecher*innen der Koalitionsparteien (SPD, Die Linke, Bündnis 90/Die Grünen) im Berliner Abgeordnetenhaus
%\item Für Hochschulen zuständiges Mitglied des Berliner Senats samt Staatssekretär
%\item Landesarbeitsgemeinschaften der Koalitionsfraktionen
%\item LandesASTenKonferenz Berlin
%\end{itemize}

%\textbf{Priorität: keine Angabe}

%\subsection*{Antrag}
%Die ZaPF möge beschließen:

\begin{quote}
Im Rahmen der geplanten Hochschulgesetz-Novellierung in Berlin hält die
Zusammenkunft aller Physikfachschaften (ZaPF) die folgenden Punkte für
besonders relevant und sieht sie als % (wichtige Grundgedanke)
Grundpfeiler eines neuen Gesetzes:


\hypertarget{studium}{%
\subsubsection*{1. Studium}\label{studium}}

\begin{enumerate}
\def\labelenumi{(\arabic{enumi})}

\item Im Zentrum des Studiums steht die Persönlichkeitsentwicklung der Studierenden. Ein Studium hat das Ziel, Menschen zum kritischen Denken und Hinterfragen anzuregen. Zudem soll der Abschluss des Studiums eine Berufsbefähigung gewährleisten. Die derzeitige \okay{alleinige} Hauptzielsetzung einer Berufsqualifizierung lehnen wir ab.

\item Alle Menschen haben ein Recht auf Bildung und freie Entfaltung. Wir lehnen daher Zugangs- und Zulassungsvoraussetzungen jeder Form ab. Die dafür notwendigen Kapazitäten an den Hochschulen sind zu schaffen.

 \emph{Für Details siehe
 }\href{https://zapfev.de/resolutionen/wise16/Zugangs-Zulassungsbeschraenkung/Reso\%20gegen\%20Zugangs-\%20und\%20Zulassungsbeschraenkungen.pdf}{\emph{Resolution
     gegen Zugangs- und Zulassungsbeschränkungen}}.\footnote{\url{https://zapfev.de/resolutionen/wise16/Zugangs-Zulassungsbeschraenkung/Reso\%20gegen\%20Zugangs-\%20und\%20Zulassungsbeschraenkungen.pdf}}
  
\item Im Sinne der freien Entfaltung und der eigenständigen Persönlichkeitsentwicklung lehnen wir jede Art der Zwangsberatung von Studierenden ab. Eine Beratung ist nur dann sinnvoll, wenn sie freiwillig und aus eigener Entscheidung erfolgt. Daher sollen Beratungsangebote beworben %werden
  und die Kapazitäten geschaffen werden.

\item Derzeit beträgt der Arbeitsaufwand%,
  bei Studienfächern, die auf die Vorlesungszeit konzentriert sind, 50 - 60 Stunden pro Woche. Dies ist die Folge einer Berechnung von Leistungspunkten an der oberen Grenze des in den Akkreditierungsrichtlinien festgelegten Bereiches von 25-30 Zeitstunden/Leistungspunkt. Für eine Normalisierung der Arbeitsbelastung soll der Maßstab von 1 LP = 25 Stunden Arbeit in Präsenz- und Selbststudium angewendet werden.

\item Momentan bemisst sich die Regelstudienzeit an einem durchschnittlichen Arbeitspensum von 30 LP pro Semester. Daraus ergibt sich bei einem Bachelor mit 180 Leistungspunkten eine Regelstudienzeit von 3 Jahren. Hierbei ist das Wort \enquote{Regel} irreführend, da aktuell die Regelstudienzeit nicht die durchschnittliche Studiendauer beschreibt, die höher liegt. Daher fordern wir die Leistungspunkte pro Semester bei 20-25 LP anzusetzen, so dass die Planstudienzeit zwischen 3.5 und 4.5 Jahre beträgt. Dies entspricht einer Annäherung an die real absolvierbaren Leistungspunkte pro Semester und erleichtert Studierenden im Vollzeitstudium, in einem Umfang arbeiten zu können, der den Lebenshaltungskosten gerecht wird.

\item Ein Teilzeitstudium ermöglicht das Studieren in Vereinbarkeit mit Arbeit, Familie und anderen Lebensbereichen. Diese Möglichkeit müssen alle Studierenden erhalten, weshalb das Recht auf ein Teilzeitstudium, ohne Begründung, im Berliner Hochschulgesetz verankert werden soll. Einzuschätzen, in welcher Zeit man sich welche Menge an Fachwissen aneignen kann, ist Teil der Persönlichkeitsentwicklung und muss in freier Entscheidung möglich sein.

\item Fachsemester definieren die Zeit, in der jemand in einen Studiengang bisher insgesamt eingeschrieben war. Bei Wechsel der Universität und Immatrikulation in den gleichen Studiengang wird die Person in das Fachsemester +1 eingestuft. Beim Wechsel des Studienganges in einen fachähnlichen oder fachnahen Studiengang wird die Person in das erste oder auf Antrag in ein höheres Fachsemester eingestuft. Die Hochschulsemester sind die Summe aller Fachsemester. Ein Teilzeitstudium führt zu einem anteiligen Fachsemester entsprechend des Verhältnisses der angestrebten Arbeitsbelastung zum Vollzeitstudium. Durch diese Abänderung der Handhabung von Fachsemestern schaffen wir Probleme beim Hochschulwechsel aus der Welt.

  \emph{Für Details siehe die
    }\href{https://zapfev.de/resolutionen/sose18/Mobilitaet/reso_mobilitaet.pdf}{\emph{Resolution zur Studierendenmobilität}}% aus Heidelberg 2018.
.\footnote{\url{https://zapfev.de/resolutionen/sose18/Mobilitaet/reso_mobilitaet.pdf}}
  \end{enumerate}

\hypertarget{semestergebuxfchren}{%
\paragraph{Semestergebühren}\label{semestergebuxfchren}}

Die ZaPF spricht sich gegen die Erhebung von Verwaltungs- und
Rückmeldegebühren aus. Wir sehen die Durchführung von
Verwaltungsaufgaben im Rahmen des Studiums wie Immatrikulation,
Prüfungsverwaltung oder Rückmeldung als Kernaufgaben der
Hochschulen, die ohne die Erhebung von Gebühren finanziert werden
müssen.


\hypertarget{pruxfcfungen}{%
\subsubsection*{2. Prüfungen}\label{pruxfcfungen}}

\hypertarget{pruxfcfungsan--und--abmeldungen}{%
\paragraph{Prüfungsan- und -abmeldungen}\label{pruxfcfungsan--und--abmeldungen}}

\begin{quote}
Prüfungsan- und abmeldungen werden von Hochschulen individuell
gehandhabt und dienen oft einem logistischen Zweck. Dies geht teilweise soweit, dass selbst innerhalb einer Hochschule oft deutliche Unterschiede zu vermerken sind. Hier stehen die Fristen im Widerspruch zu Flexibilität und Studierendenfreundlichkeit.

In unseren Augen gibt es keinen Grund, warum Studierende zum Teil
mehrere Wochen vor Prüfungstermin von einer Prüfungsanmeldung
zurücktreten müssen und wir sehen in dieser Form der Handhabung unnötige Hürden für Studierende. Eine Prüfungsanmeldung soll, falls sie denn explizit nötig ist, revidierbar sein. Diese Revision sollte so spät wie möglich vor der Prüfung durchführbar sein.
\end{quote}

\emph{Für ausführliche Details zu den Überlegungen der ZaPF siehe die}
\href{https://zapfev.de/resolutionen/sose18/Pruefungsanmeldung/reso_pruefungsanmeldung.pdf}{\emph{Resolution
    für einen flexibleren Umgang mit Prüfungsan- und abmeldungen}}.\footnote{\url{https://zapfev.de/resolutionen/sose18/Pruefungsanmeldung/reso_pruefungsanmeldung.pdf}}

\hypertarget{zwangsexmatrikulation}{%
\paragraph{Zwangsexmatrikulation}\label{zwangsexmatrikulation}}

\begin{quote}
Studierende durch drohende Zwangsexmatrikulation unter Druck zu setzen ist in unseren Augen unangemessen; es ersetzt selbstverantwortliches und selbstbestimmtes durch prüfungsorientiertes Studieren und behindert damit die freie Entfaltung. Die ZaPF spricht sich gegen sämtliche Regelungen in Studienordnungen aus, welche den Fokus des Studiums von der Aneignung von Wissen und persönlicher Entwicklung hin zu der
Verhinderung der eigenen Exmatrikulation verschieben. Insbesondere fordern wir, solche Regelungen aufzuheben oder abzuändern, die auf eine Zwangsexmatrikulation hinaus laufen können, insbesondere die Begrenzung der Anzahl von Prüfungsversuchen oder die Zwangsberatungen.
\end{quote}

\emph{Ausführliche Details zu den Überlegungen der ZaPF findet man in der} \href{https://zapfev.de/resolutionen/wise17/Zwangsexmatrikulation/Zwangsexmatrikulation.pdf}{\emph{Resolution
zur Zwangsexmatrikulation}}.\footnote{\url{https://zapfev.de/resolutionen/wise17/Zwangsexmatrikulation/Zwangsexmatrikulation.pdf}}

\hypertarget{symptompflicht}{%
\paragraph{Symptompflicht}\label{symptompflicht}}

\begin{quote}
Die ZaPF spricht sich gegen die geforderte Angabe von Symptomen auf Attesten für die Prüfungsunfähigkeitsmeldung aus. Wir fordern, ausschließlich das Verfahren der Prüfungsunfähigkeitsbescheinigungen analog zu Arbeitsunfähigkeitsbescheinigungen zu ermöglichen.
\end{quote}

\emph{Für ausführliche Details zu den Überlegungen der ZaPF siehe das} \href{https://zapfev.de/resolutionen/sose17/symptompflicht/PosPapier_Symptompflicht.pdf}{\emph{Positionspapier
zur Symptompflicht auf Attesten}}\footnote{\url{https://zapfev.de/resolutionen/sose17/symptompflicht/PosPapier_Symptompflicht.pdf}} \emph{und die} \href{https://zapfev.de/resolutionen/wise17/Pruefungsunfaehigkeit/Pruefungsunfaehigkeitsbescheinigung.pdf}{\emph{Resolution
zu Prüfungsunfähigkeitsbescheinigungen}}.\footnote{\url{https://zapfev.de/resolutionen/wise17/Pruefungsunfaehigkeit/Pruefungsunfaehigkeitsbescheinigung.pdf}}


\hypertarget{akkreditierung}{%
\subsubsection*{3. Akkreditierung}\label{akkreditierung}}

Die Qualitätssicherung im Rahmen von (System-)Akkreditierungen muss eine Rückbindung an die demokratischen Hochschulgremien erfahren, da sonst Verwaltungsstellen über die Ausgestaltung und Qualitätssicherung von Studiengängen entscheiden. Die Einbindung der Ausbildungskommissionen und der Kommissionen für Studium und Lehre muss gesetzlich garantiert werden.

Neuakkreditierungen sollen mit einer verkürzten Frist von fünf Jahren
gelten. \emph{Für mehr Details siehe die }\href{https://zapfev.de/resolutionen/sose18/Akkreditierung/reso_laender_akkr.pdf}{\emph{Resolution zur länderspezifischen Ausgestaltung der MRVO}}.\footnote{\url{https://zapfev.de/resolutionen/sose18/Akkreditierung/reso_laender_akkr.pdf}} Alle Gutachter*innen sollen im Bereich Akkreditierung geschult sein -- entweder durch ihre Erfahrung oder durch entsprechende Fortbildungsmaßnahmen. Bei Akkreditierungen von Lehramtsstudiengängen darf die Vertretung der Berufspraxis in der Gutachtergruppe nicht durch einen Vertreter*in der obersten Landesbehörde ersetzt werden, sondern soll um diese*n ergänzt werden.


\hypertarget{verfasste-studierendenschaft}{%
\subsubsection*{4. Verfasste Studierendenschaft}\label{verfasste-studierendenschaft}}

\hypertarget{stuxe4rkung-des-zugangs-zu-hochschuleinrichtungen}{%
\paragraph{Stärkung des Zugangs zu Hochschuleinrichtungen}\label{stuxe4rkung-des-zugangs-zu-hochschuleinrichtungen}}

Verfasste Studierendenschaften müssen für die Erfüllung ihrer
vielfältigen und wichtigen Aufgaben freien Zugang zu den Einrichtungen der Hochschule haben. Insbesondere die ungehinderte Nutzung von Räumen für Veranstaltungen und der Versand von Mails an die Studierenden sind Rechte, die im Gesetz garantiert werden müssen. Die Hochschule hat die verfasste Studierendenschaft durch Zugang zu allgemeinen Verwaltungseinrichtungen wie Briefversand z.B. an neue Erstsemester zu unterstützen.

\hypertarget{fachschaftsinitativen}{%
\paragraph{Fachschaftsinitativen}\label{fachschaftsinitativen}}

Viele Hochschulen in Berlin haben als innovatives Konzept das offene Modell der studentischen Partizipation eingeführt, die sogenannten Fachschaftsinitativen. Diese niederschwelligen Strukturen sind als studentische Gestaltungsmöglichkeiten in der akademischen Selbstverwaltung entsprechend im Gesetz aufzunehmen. Ihnen muss der Zugang zu den Einrichtungen der Hochschulen garantiert werden, insbesondere die Bereitstellung von Fachschaftsräumen.

\hypertarget{semesterticket}{%
\paragraph{Semesterticket}\label{semesterticket}}

Zugang zum öffentlichen Personennahverkehr ist eine unverzichtbare Voraussetzung für ein sozialverträgliches Studium. In Anbetracht der steigenden Lebenshaltungskosten in Berlin sehen wir es als notwendig, dass ein Semesterticket nicht nur kostengünstig, sondern höchstens kostenneutral ist.


\hypertarget{gremien}{%
\subsubsection*{5. Gremien}\label{gremien}}

\hypertarget{zusammensetzung-und-statusgruppenveto}{%
\paragraph{Zusammensetzung und Statusgruppenveto}\label{zusammensetzung-und-statusgruppenveto}}

% Im Sinne einer demokratischen Hochschule, sind folgende Punkte wichtig zu beachten:

Für die konstruktive Zusammenarbeit aller Statusgruppen in Universitätsgremien empfiehlt die ZaPF folgende Punkte zu beachten:

\begin{quote}
\begin{enumerate}
\def\labelenumi{\arabic{enumi}.}
%\tightlist
\item Grundsätzlich ist es falsch, wenn eine Statusgruppe in einem demokratischen Gremium automatisch die Mehrheit besitzt. Vielmehr ist es notwendig, dass keine Position übergangen werden kann. Dies kann z.B. durch eine paritätische Zusammensetzung oder ein Statusgruppen-Vetorecht {[}(siehe Punkt 2){]} sicher gestellt werden. Voraussetzung dafür ist, dass die Teilhaberechte aller gesetzlich sichergestellt sind und nicht nur optional gewährt werden.
\end{enumerate}
\end{quote}

\emph{Vergleiche hierzu }\href{https://zapfev.de/resolutionen/sose18/Hochschulgesetze/reso_hsgesetze.pdf}{\emph{Abschnitt \enquote{Gremien} der Resolution zu Hochschulgesetzen%, verabschiedet in Heidelberg 2018
  }}.\footnote{\url{https://zapfev.de/resolutionen/sose18/Hochschulgesetze/reso_hsgesetze.pdf}}

\begin{quote}
\begin{enumerate}
\def\labelenumi{\arabic{enumi}.}
\setcounter{enumi}{1}
%\tightlist
\item Lehnt eine Statusgruppe geschlossen einen Antrag ab, soll sie das Recht haben, ein Veto einzulegen (\emph{Statusgruppenveto}). %{[}\ldots{}{]}

\item Die Mitglieder des Dekanats {[}und des Präsidiums bzw. Rektorats{]} dürfen ausschließlich mit beratender Funktion an Gremiumssitzungen teilnehmen. {[}Die Leitung von Gremien erfolgt durch einen selbst gewählten Vorsitz.{]}
\end{enumerate}
\end{quote}

\href{https://zapf.wiki/Sammlung_aller_Resolutionen_und_Positionspapiere\#Positionspapier_zur_demokratischen_Mitgestaltung_in_Hochschulgremien}{\emph{Vergleiche das Positionspapier zur demokratischen Mitgestaltung in Hochschulgremien%aus Jena 2013
  }}.\footnote{\url{https://zapf.wiki/Sammlung_aller_Resolutionen_und_Positionspapiere\#Positionspapier_zur_demokratischen_Mitgestaltung_in_Hochschulgremien}}

\begin{quote}
\begin{enumerate}
\def\labelenumi{\arabic{enumi}.}
\setcounter{enumi}{3}
%\tightlist
\item Gewählten Stellvertretungen der Mitglieder aller Gremien darf die Anwesenheit auch in nichtöffentlichen Sitzungen des jeweiligen Gremiums nicht verwehrt werden.
\end{enumerate}
\end{quote}

\hypertarget{kontrollrechte-von-gremienmitgliedern}{%
\paragraph{Kontrollrechte von Gremienmitgliedern}\label{kontrollrechte-von-gremienmitgliedern}}

Zur Wahrung der demokratischen Grundsätze in der akademischen Selbstverwaltung sind Kontrollrechte für Gremienmitglieder unabdingbar. Daher fordern wir das Recht für Gremienmitglieder, Berichte sowie Akteneinsicht zu allen in den Zuständigkeitsbereich des Gremiums und der von ihm gewählten Ämter fallenden Fragestellungen zu erhalten. Den Gremien muss darüber hinaus das Recht auf Einholung von Gutachten und Stellungnahmen garantiert werden.

\hypertarget{gremienvor--beziehungsweise--nachmittag}{%
\paragraph{Gremienvor- beziehungsweise -nachmittag}\label{gremienvor--beziehungsweise--nachmittag}}

Die ZaPF spricht sich dafür aus, im Hochschulgesetz einen für
Gremienarbeit reservierten Vor- oder Nachmittag zu verankern.
Insbesondere sollen in diesem Zeitraum keine Lehrveranstaltungen
stattfinden. Die Wahl des Zeitraums und des Wochentages soll der
jeweiligen Hochschule überlassen bleiben.


\hypertarget{gleichstellung}{%
\subsubsection*{6. Gleichstellung}\label{gleichstellung}}

Die Forderungen der ZaPF zur Durchsetzung von Gleichstellung an den Berliner
Hochschulen lauten:
\begin{enumerate}
\item[(1)] genderneutrale Formulierungen im BerlHG. Beispiel:
  \enquote{Studierende} statt \enquote{Studentinnen und Studenten}.
\item[(2)] dass Hochschulgrade in genderneutraler Form verliehen werden.
\item[(3)] eine Änderung der Regelung zur Frauenbeauftragten, so dass die Anzahl
  der hauptberuflichen Frauenbeauftragten von der Anzahl der Studierenden der
  Hochschule abhängt, mindestens jedoch drei hauptberufliche Frauenbeauftragte
  pro Hochschule tätig sind.
\item[(4)] zur Wahl der haupt- und nebenberuflichen Frauenbeauftragten haben
  alle Mitglieder der Hochschule das passive Wahlrecht. Das aktive Wahlrecht
  setzt den Eintrag \glqq weiblich'' in der Geburtsurkunde voraus.
\item[(5)] eine hauptberuflich beauftragte Person zur Gender-Gleichstellung.
\item[(6)] eine beauftragte Person für Studierende mit Behinderungen %nach BerlHG
  pro Fachbereich.
\item[(7)] eine beauftragte Person für Studierende mit Kindern, analog zu (5)
  pro Fachbereich.
\item[(8)] Eltern-Kind-Zimmer, die KiTa-Standards erfüllen.
\item[(9)] die Möglichkeit, bei der Immatrikulation nicht-binäre Genderoptionen
  wählen zu können.
\end{enumerate}

\hypertarget{promotion}{%
\subsubsection*{7. Promotion}\label{promotion}}

\hypertarget{rolle-und-statusgruppe}{%
\paragraph{Rolle und Statusgruppe}\label{rolle-und-statusgruppe}}

Die Promotion dient dem Nachweis der Befähigung zu vertiefter wissenschaftlicher Arbeit. Sie ist die erste Phase selbständiger wissenschaftlicher Arbeit. 

Wir als ZaPF fordern die Aufnahme dieser Klarstellung in das Berliner Hochschulgesetz und in Konsequenz die Einordnung aller Promovierender in die Statusgruppe der wissenschaftlich %en
Mitarbeitenden.

\hypertarget{bessere-promotionsbedingungen}{%
\paragraph{Bessere Promotionsbedingungen}\label{bessere-promotionsbedingungen}}

Zur Auflösung der Abhängigkeit von der betreuenden Person fordern wir die personelle Trennung von Betreuung, Begutachtung und arbeitsrechtlicher Weisungsbefugnis.

Um dies zu gewährleisten müssen Promovierendenzentren an den Hochschulen mit den Aufgaben:
\begin{itemize}
\item Zulassung zur Promotion und Zuordnung der fachlichen Betreuung,
\item Finanzielle und arbeitsrechtliche Organisation des Promotionsvorhabens,
\item Bestellung der Gutachter der Doktorarbeit
\item Unterstützung und Weiterbildung in Fragen der Lehre, 
\item Überfachliche Weiterbildung,
\item Ombudsstelle
\end{itemize}
eingerichtet werden.

Die Lehrtätigkeit Promovierender ist eine der Säulen universitärer Lehre. Die Qualität dieser Lehrtätigkeit ist daher von besonderer Wichtigkeit. Aus diesem Grund müssen Promovierende bei der Entwicklung ihrer didaktischen Fähigkeiten unterstützt werden.

Als Anstellungsverhältnis fordern wir Qualifizierungsstellen mit den unter Arbeitsbedingungen genannten Standards (siehe \href{https://protokolle.zapf.in/reso_BerlHG\#Qualifizierungsstellen}{unten}).


\hypertarget{arbeitsbedingungen-in-der-wissenschaft}{%
\subsubsection*{8. Arbeitsbedingungen in der Wissenschaft}\label{arbeitsbedingungen-in-der-wissenschaft}}

\hypertarget{qualifizierungsstellen}{%
\paragraph{Qualifizierungsstellen}\label{qualifizierungsstellen}}

%\emph{Arbeitsverhältnisse auf Qualifizierungsstellen an Hochschulen werden grundsätzlich auf Basis von 100 Prozent der regulären tariflichen Arbeitszeit geschlossen. Begründete Ausnahmen hiervon (Teilzeit) bedürfen der Zustimmung des Arbeitnehmers bzw. der Arbeitnehmerin. Wird ein Arbeitsverhältnis mit weniger als 100 Prozent der regulären tariflichen Arbeitszeit geschlossen, so ist eine Mehrarbeit über das vertraglich festgesetzte Maß unzulässig. Es stehen mindestens 50 Prozent der regulären tariflichen Arbeitszeit zur Erreichung des Qualifikationszieles zur Verfügung.}

Für Arbeitsverhältnisse mit einem Qualifikationsziel fordern wir folgende Standards:

Arbeitsverhältnisse auf Qualifizierungsstellen an Hochschulen werden grundsätzlich auf Basis von 100 Prozent der regulären tariflichen Arbeitszeit geschlossen. In begründeten Ausnahmen (Teilzeit) darf hiervon in Absprache mit den Arbeitnehmenden abgewichen werden. Wird ein Arbeitsverhältnis mit weniger als 100 Prozent der regulären tariflichen Arbeitszeit geschlossen, so ist eine Mehrarbeit über das vertraglich festgesetzte Maß unzulässig. Es stehen mindestens 50 Prozent der regulären tariflichen Arbeitszeit zur Erreichung des Qualifikationszieles zur Verfügung.

\hypertarget{dauerstellen}{%
\paragraph{Dauerstellen}\label{dauerstellen}}
Die ZaPF fordert die Schaffung unbefristeter Stellen im wissenschaftlichen Mittelbau.

\begin{quote}
Nur durch eine deutliche Erhöhung der Anzahl unbefristeter Stellen im wissenschaftlichen Mittelbau kann es zu einer nachhaltigen Qualitätssicherung in der Forschung und Lehre, effizientem Wissenstransfer und einer Steigerung der Attraktivität der Karriere in der Wissenschaft kommen.
\end{quote}

\href{https://zapfev.de/resolutionen/sose17/mittelbau/mittelbau.pdf}{\emph{Siehe Resolution zur Schaffung permanenter
    Stellen im wissenschaftlichen Mittelbau%, die am 27.05.2017 in Berlin verabschiedet wurde.
  }}\footnote{\url{https://zapfev.de/resolutionen/sose17/mittelbau/mittelbau.pdf}}

\begin{quote}
Insbesondere müssen Daueraufgaben durch unbefristete Anstellungen abgedeckt sein.
\end{quote}

\href{https://zapfev.de/resolutionen/wise15/WissZeitVG/Stellungnahme_WiSe15_WissZeitVG.pdf}{\emph{Siehe Resolution zur Novelle des Wissenschaftszeitvertragsgesetzes%, beschlossen in Frankfurt 2015 (Z. 30-35)
  }}.\footnote{\url{https://zapfev.de/resolutionen/wise15/WissZeitVG/Stellungnahme_WiSe15_WissZeitVG.pdf}}

\hypertarget{gute-wissenschaftliche-praxis}{%
\paragraph{Gute Wissenschaftliche Praxis}\label{gute-wissenschaftliche-praxis}}
Forschung nach den Prinzipien guter wissenschaftlicher Praxis ist in allen Bereichen der Wissenschaft anzustreben.

Als konkrete Maßnahme sehen wir die Notwendigkeit der Einrichtung unabhängiger Ombudsstellen auf Landesebene.


\hypertarget{transparenz-und-gesellschaftliche-verantwortung-der-hochschulen}{%
\subsubsection*{9. Transparenz und Gesellschaftliche Verantwortung der Hochschulen}\label{transparenz-und-gesellschaftliche-verantwortung-der-hochschulen}}

\hypertarget{transparenz-in-der-drittmittel-finanzierten-forschung}{%
\paragraph{Transparenz in der Drittmittel-finanzierten Forschung}\label{transparenz-in-der-drittmittel-finanzierten-forschung}}

\begin{quote}
Wir halten Transparenz bei der Durchführung von wissenschaftlichen Tätigkeiten im Interesse Dritter für notwendig. Deshalb fordert die ZaPF, dass bei Drittmittelprojekten folgende Angaben jährlich veröffentlicht werden müssen:
\begin{enumerate}
\item Titel des Projekts
\item Hochschule mit Organisationseinheit
\item Auftraggebende Personen mit Sparte/Handlungsfeld der Abteilung
\item Projekt- und Vertragslaufzeit
\item Gesamtsumme
\item Angaben der Geheimhaltungsvereinbarungen oder Publikationsbeschränkungen,
  u. a. Art, Dauer und Umfang
\end{enumerate}
\end{quote}

\href{https://zapfev.de/resolutionen/wise15/Transparenz_in_der_Drittmittelforschung/Stellungnahme_WiSe15_Transparenz_in_der_Drittmittelforschung.pdf}{\emph{Positionspapier zur Transparenz in der Drittmittelforschung% verabschiedet am 22.11.2015 in Frankfurt% mit Streichungen
  }}\footnote{\url{https://zapfev.de/resolutionen/wise15/Transparenz_in_der_Drittmittelforschung/Stellungnahme_WiSe15_Transparenz_in_der_Drittmittelforschung.pdf}}

Zusätzlich muss am Projektende ein Abschlussbericht veröffentlicht werden.

\hypertarget{friedensbindung}{%
\paragraph{Friedensbindung}\label{friedensbindung}}
Im Berliner Hochschulgesetz fehlt bisher die Verpflichtung zur zivilen Forschung. Diese Verantwortung soll bei den Aufgaben der Hochschule im Gesetz verankert werden.


\hypertarget{psychologische-erstberatung}{%
\subsubsection*{10. Psychologische Erstberatung}\label{psychologische-erstberatung}}

Mit dem Studium beginnt ein neuer Lebensabschnitt in dem Menschen mit neuen Herausforderungen konfrontiert werden. Dies kann zu psychischen Belastungen führen, die unter Umständen nicht ohne professionelle Ansprechpersonen bewältigt werden können. Für diese Aufgabe muss die Hochschule eine psychologische Anlauf- und Beratungsstelle einrichten. Zudem sieht die ZaPF die Sensibilisierung im Umgang mit psychologischen Problemen als wichtiges gesellschaftliches Thema und unterstützt die Schaffung und aktive Bewerbung von universitären Beratungsstellen.
\end{quote}

\end{document}

%%% Local Variables:
%%% mode: latex
%%% TeX-master: t
%%% End:
